\chapter{اعداد}
\thispagestyle{empty}

\begin{quote}
	{\lotus
		اعداد، مخلوقات آزاد ذهن بشر هستند. \quad - ددکیند%
		\LTRfootnote{ Julius Wilhelm Richard Dedekind}
	}
\end{quote}

\section{مجموعه‌های اعداد}
مجموعه‌ی اعداد طبیعی%
\RTLfootnote{ در این جزوه، صفر را جزو اعداد طبیعی می‌دانیم. درباره‌ی سایر منابع به قرارداد خود آن منابع مراجعه کنید.}:
$$\mathbb{N} = \{0,1,2,3, \cdots\}$$
مجموعه‌ی اعداد صحیح:
$$\mathbb{Z} = \{\cdots, -3,-2,-1,0,1,2,3, \cdots\}$$
مجموعه‌ی اعداد صحیح مثبت:
$$\mathbb{Z}^+ = \{1,2,3, \cdots\}$$
مجموعه‌ی اعداد صحیح منفی:
$$\mathbb{Z}^- = \{\cdots, -3,-2,-1\}$$
مجموعه‌ی اعداد گویا:
$$\mathbb{Q} = \left\{\frac{m}{n} \ | \ m,n \in \mathbb{Z} \ \wedge \ n \neq 0\right\}$$

\section{اصول موضوعه‌ی اعداد حقیقی}
اعداد حقیقی را با نماد
$\mathbb{R}$
نشان می‌دهیم. در اینجا اعداد حقیقی را تعریف نمی‌کنیم، بلکه اصول موضوعه‌ای را برای آن بیان می‌کنیم.
\begin{enumerate}
	\item
	بسته‌بودن مجموعه‌ی
	$\mathbb{R}$
	نسبت به عمل جمع:
	$$\forall a,b \in \mathbb{R} ,\quad a+b \in \mathbb{R}$$
	
	\item
	بسته‌بودن مجموعه‌ی
	$\mathbb{R}$
	نسبت به عمل ضرب:
	$$\forall a,b \in \mathbb{R} ,\quad a \times b \in \mathbb{R}$$
	
	\item 
	جابه‌جایی عمل جمع:
	$$\forall a,b \in \mathbb{R} ,\quad a+b = b+a$$
	
	\item 
	جابه‌جایی عمل ضرب:
	$$\forall a,b \in \mathbb{R} ,\quad a \times b = b \times a$$
	
	\item 
	شرکت‌پذیری عمل جمع:
	$$\forall a,b,c \in \mathbb{R} ,\quad (a+b)+c = a+(b+c)$$
	
	\item 
	شرکت‌پذیری عمل ضرب:
	$$\forall a,b,c \in \mathbb{R} ,\quad (a \times b) \times c = a \times (b \times c)$$
	
	\item 
	عضو خنثای عمل جمع: فقط یک عدد حقیقی وجود دارد که برای عمل جمع اعداد حقیقی، خنثی است. این عدد را صفر می‌نامیم و آن را با نماد $0$ نشان می‌دهیم.
	$$\forall a \in \mathbb{R} ,\quad a+0 = 0+a = a$$
	
	\item 
	عضو خنثای عمل ضرب: فقط یک عدد حقیقی مخالف با صفر وجود دارد که برای عمل ضرب اعداد حقیقی، خنثی است. این عدد را یک می‌نامیم و آن را با نماد $1$ نشان می‌دهیم.
	$$\forall a \in \mathbb{R} ,\quad a \times 1 = 1 \times a = a$$
	
	\item 
	وجود عدد قرینه: برای هر عدد $a$ فقط یک عدد وجود دارد که مجموع آن با $a$ مساوی است با صفر. این عدد را قرینه‌ی عدد $a$ می‌نامیم و آن را با نماد $-a$ نشان می‌دهیم.
	$$\forall a \in \mathbb{R} ,\quad a+(-a) = (-a)+a = 0$$
	
	\item 
	وجود وارون: برای هر عدد حقیقی $a$ مخالف صفر فقط یک عدد حقیقی وجود دارد که حاصل‌ضرب آن در $a$ مساوی است با یک. این عدد را وارون عدد $a$ می‌نامیم و آن را با نماد
	$\frac{1}{a}$ یا $a^{-1}$
	نشان می‌دهیم.
	$$\forall a \in \mathbb{R}-\{0\} ,\quad a \times (\frac{1}{a}) = (\frac{1}{a}) \times a = 1$$
	
	\item 
	خاصیت پخشی (توزیع‌پذیری) عمل ضرب نسبت به عمل جمع:
	$$\forall a,b,c \in \mathbb{R} ,\quad a \times (b+c) = (a \times b) + (a \times c)$$
	$$\forall a,b,c \in \mathbb{R} ,\quad (b+c) \times a = (b \times a) + (c \times a)$$
\end{enumerate}

\begin{theorem}
	قرینه‌ی قرینه‌ی هر عدد حقیقی مساوی است با خود آن عدد، یعنی:
	$-(-a)=a$.
	\\
	\textbf{برهان:}
	بنا به اصل ۹ می‌دانیم قرینه‌ی $a$ برابر است با $-a$، یعنی مجموع $a$ و $-a$ برابر صفر است. از طرفی قرینه‌ی $-a$ یعنی	$-(-a)$ یکتاست، پس فقط یک عدد وجود دارد که جمع آن با $-a$ برابر صفر است. از قبل می‌دانیم که مجموع $a$ و $-a$ برابر صفر است. پس قرینه‌ی $-a$ برابر با $a$ است.
	\qed
\end{theorem}

\begin{theorem}
	قرینه‌ی مجموع دو عدد حقیقی مساوی است با مجموع قرینه‌های آن‌ها، یعنی:
	$-(a+b)=(-a)+(-b)$.
	\\
	\textbf{برهان:}
	بنا به اصل ۹ داریم:
	$$(a+b)+(-(a+b))=0$$
	همچنین داریم:
	\begin{align*}
		(a+b)+((-a)+(-b)) &= a+b+(-a)+(-b) \\
		&= a+(-a)+b+(-b) &\text{اصل ۳} \\
		&= (a+(-a))+(b+(-b)) &\text{اصل ۵} \\
		&= 0+0 &\text{اصل ۹} \\
		&= 0
	\end{align*}
	با توجه به این‌که عضو قرینه یکتاست، بنابراین:
	$-(a+b)=(-a)+(-b)$.
	\qed
	\\ \df \\
	\textbf{نتیجه:}
	برای تعدادی متناهی جمله داریم (سه نقطه تا ابد ادامه ندارد):
	$$-(a+b+c+\cdots) = (-a)+(-b)+(-c)+\cdots$$
\end{theorem}
\smallskip
\textbf{عمل تفریق (تفاضل):}
اگر $a$ و $b$ دو عدد حقیقی باشند، مجموع عدد $a$ و قرینه‌ی عدد $b$ یعنی $-b$ را تفریق (تفاضل) $a$ و $b$ می‌نامند و آن را با
$a-b$
نمایش می‌دهند. تفریق خاصیت جابه‌جایی ندارد و ترتیب در آن مهم است.
\\
\textbf{عمل تقسیم:}
اگر $a$ و $b$ دو عدد حقیقی باشند و
$b \neq 0$،
حاصل‌ضرب عدد $a$ در وارون عدد $b$ یعنی
$\frac{1}{b}$
را تقسیم $a$ بر $b$ می‌نامند و آن را با
$\frac{a}{b}$ یا $a \div b$
نمایش می‌دهند. تقسیم خاصیت جابه‌جایی ندارد و ترتیب در آن مهم است.

\begin{theorem}
	حاصل‌ضرب دو عدد حقیقی برابر با صفر است اگر و تنها اگر دست‌کم یکی از آن دو صفر باشد.
	$$\forall a,b \in \mathbb{R} ,\quad (a=0) \vee (b=0) \iff a \times b=0$$
\end{theorem}
